%------------ Debut Packages ----------------
\usepackage{lmodern}
\usepackage[table]{xcolor} % pour ajouter de la couleur (si besoin)
\usepackage{multirow}
\usepackage{makecell}
\usepackage{float}
\usepackage{tablefootnote}
\usepackage{bookmark}
\usepackage{pdfpages}
\usepackage[linesnumbered,ruled,vlined]{algorithm2e}
\usepackage{lipsum}
\usepackage{listings}
\usepackage{lscape}
\usepackage{url}
\usepackage{amsmath}
\usepackage{amssymb}
\usepackage{fontawesome5}
\usepackage{tabularx}
\usepackage{ragged2e}
\usepackage[most]{tcolorbox}

% --- Custom Colors ---
\definecolor{Navy}{RGB}{15,23,42}
\definecolor{Teal}{RGB}{20,184,166}
\definecolor{Slate}{RGB}{71,85,105}
\definecolor{LightGray}{RGB}{248,250,252}
\definecolor{MediumGray}{RGB}{230,235,240}
\definecolor{CodeBg}{RGB}{241,245,249}
\definecolor{EMSIGreen}{RGB}{12, 114, 56} 

% --- Tabularx Column Types ---
\newcolumntype{L}{>{\bfseries\RaggedRight\arraybackslash}p{3.6cm}}
\newcolumntype{Y}{>{\RaggedRight\arraybackslash}X}

% --- TikZ Libraries ---
\usepackage{tikz}
\usetikzlibrary{
    shapes.geometric,
    shapes.symbols,
    shapes.multipart,
    arrows.meta,
    positioning,
    fit,
    backgrounds,
    shadows,
    calc,
    matrix
}

% --- Use Case Box ---
\newtcolorbox{usecasebox}[1]{
    enhanced,
    colback=white,
    colframe=black,
    boxrule=0.6pt,
    arc=2mm,
    left=3mm,
    right=3mm,
    top=4mm,
    bottom=3mm,
    title={#1},
    fonttitle=\bfseries,
    coltitle=black,
    colbacktitle=MediumGray,
    attach boxed title to top left={
        xshift=3mm,
        yshift=-2mm
    },
    boxed title style={
        colframe=black,
        colback=MediumGray,
        boxrule=0.6pt,
        arc=1mm
    }
}

% --- Listings Settings ---
\lstset{
    backgroundcolor=\color{CodeBg},
    basicstyle=\ttfamily\footnotesize\color{Navy},
    keywordstyle=\color{Teal}\bfseries,
    stringstyle=\color{Slate},
    commentstyle=\color{gray}\itshape,
    numbers=left,
    numberstyle=\tiny\color{gray},
    breaklines=true,
    frame=single,
    rulecolor=\color{Navy!20},
    captionpos=b,
    showstringspaces=false,
    tabsize=2
}

% cleveref should be loaded after hyperref (loaded in .cls)
\usepackage[nameinlink]{cleveref}

% On définit le style pour les pages "normales"
\pagestyle{fancy}
\renewcommand{\chaptermark}[1]{\markboth{\chaptername \ \thechapter.\ #1}{}} % sert à personnaliser l'affichage de \leftmark (ici : le mot "Chapitre", le numéro, un point, et le titre du chapitre, sans écrire en majuscules)
\renewcommand{\sectionmark}[1]{\markright{\thesection.\ #1}} % sert à personnaliser l'affichage de \rightmark (ici : le numéro et le titre de la section en cours, sans écrire en majuscules),chaptermark sert à dire :Quel titre de chapitre doit apparaître en haut de la page ?
%\renewcommand{\thesection}{\arabic{section}}
\usepackage{chngcntr}
\counterwithin{section}{chapter}

%------------ Fin Packages ----------------